% Mechanistic Validity — NEMI 2026 poster, Option B
% Results-first framing: "36 Criteria, 16 Claims, No Claim Reaches Validated"
% 3-column, 3-row. Text and tables only — no figures.
%
%   cd poster && lualatex mechval_poster_option_b.tex

\documentclass[25pt, a0paper, portrait]{tikzposter}

\usepackage{amsmath,amssymb}
\usepackage{array}
\usepackage{booktabs}
\usepackage{enumitem}
\usepackage{microtype}

% ── Color palette ───────────────────────────────────────
\definecolor{confirmed}{HTML}{27AE60}
\definecolor{partial}{HTML}{2980B9}
\definecolor{inconclusive}{HTML}{F39C12}
\definecolor{untested}{HTML}{95A5A6}
\definecolor{disconfirmed}{HTML}{E74C3C}
\definecolor{darktext}{HTML}{2C3E50}
\definecolor{blockbg}{HTML}{FFFFFF}
\definecolor{tierProposed}{HTML}{95A5A6}
\definecolor{tierCausal}{HTML}{F39C12}
\definecolor{tierSupported}{HTML}{2980B9}
\definecolor{tierTriangulated}{HTML}{27AE60}
\definecolor{tierValidated}{HTML}{8E44AD}
\definecolor{tierDisconfirmed}{HTML}{E74C3C}

% ── Theme ───────────────────────────────────────────────
\usetheme{Default}

\definecolorstyle{MechValColors}{
  \definecolor{colorOne}{HTML}{2C3E50}
  \definecolor{colorTwo}{HTML}{3498DB}
  \definecolor{colorThree}{HTML}{ECF0F1}
}{
  \colorlet{backgroundcolor}{colorThree}
  \colorlet{framecolor}{colorOne}
  \colorlet{titlefgcolor}{white}
  \colorlet{titlebgcolor}{colorOne}
  \colorlet{blocktitlebgcolor}{colorTwo}
  \colorlet{blocktitlefgcolor}{white}
  \colorlet{blockbodybgcolor}{blockbg}
  \colorlet{blockbodyfgcolor}{darktext}
  \colorlet{notefgcolor}{darktext}
  \colorlet{notebgcolor}{inconclusive!15}
  \colorlet{notefrcolor}{inconclusive}
}
\usecolorstyle{MechValColors}

\defineblockstyle{MechValBlock}{
  titlewidthscale=1.0, bodywidthscale=1.0, titleleft,
  titleoffsetx=0pt, titleoffsety=0pt,
  bodyoffsetx=0pt, bodyoffsety=0pt, bodyverticalshift=0pt,
  roundedcorners=12, linewidth=2pt, titleinnersep=10pt, bodyinnersep=15pt
}{
  \draw[rounded corners=\blockroundedcorners, color=blocktitlebgcolor,
        fill=blockbodybgcolor, line width=\blocklinewidth]
    (blockbody.south west) rectangle (blocktitle.north east);
  \ifBlockHasTitle
    \draw[rounded corners=\blockroundedcorners,
          fill=blocktitlebgcolor, color=blocktitlebgcolor,
          line width=\blocklinewidth]
      (blocktitle.south west) rectangle (blocktitle.north east);
  \fi
}
\useblockstyle{MechValBlock}

\definetitlestyle{MechValTitle}{
  width=\paperwidth, roundedcorners=0, linewidth=0pt, innersep=20pt,
  titletotopverticalspace=15mm, titletoblockverticalspace=25mm
}{
  \draw[fill=titlebgcolor, color=titlebgcolor]
    (\titleposleft, \titleposbottom)
    rectangle (\titleposright, \titlepostop);
}
\usetitlestyle{MechValTitle}

\tikzposterlatexaffectionproofoff

\title{%
  \parbox{0.92\linewidth}{\centering\bfseries
    Mechanistic Validity\\[6pt]
    36 Criteria, 16 Claims, No Claim Reaches Validated}}
\author{Elliot Tower\quad\texttt{elliot@elliottower.ai}}
\institute{}

\begin{document}

\maketitle[width=\paperwidth]

% ════════════════════════════════════════════════════════
% ROW 1 — FRAMEWORK SUMMARY
% ════════════════════════════════════════════════════════
\begin{columns}
\column{0.33}

\block{Five Validity Types}{
  Every mechanistic claim owes evidence along five dimensions,
  each imported from an established field:

  \vspace{4mm}
  \renewcommand{\arraystretch}{1.25}
  \begin{tabular}{@{}lp{0.62\linewidth}@{}}
    \toprule
    \textbf{Type} & \textbf{Core question} \\
    \midrule
    Construct    & Is the circuit a natural kind? \\
    Measurement  & Does the metric measure what it claims? \\
    Internal     & Is the causal claim warranted? \\
    External     & Does the result generalize? \\
    Interpretive & Is the label justified? \\
    \bottomrule
  \end{tabular}

  \vspace{5mm}
  Construct and measurement validity come from Cronbach \&
  Meehl. Internal validity from Campbell. External from
  Shadish et al.

  \vspace{4mm}
  \textbf{Interpretive validity is new.} A claim about how a
  model works can be well measured, causally established, and
  still wrong about what its components do---because the label
  asserts more than the intervention that identified the
  component can support.

  \vspace{4mm}
  The five types form a \textbf{dependency chain}: construct
  validity is a ceiling on all the rest.
}

\column{0.34}

{
\colorlet{blocktitlebgcolor}{tierDisconfirmed}
\block{The Minimum Rule}{
  {\Large\bfseries A verdict is the weakest validity dimension,
  not the average.}

  \vspace{6mm}
  A single strong dimension cannot compensate for an untested
  one. Strong internal validity with untested construct validity
  means the causal result is real but the target concept may be
  incoherent.

  \vspace{6mm}
  \begin{center}
  \renewcommand{\arraystretch}{1.35}
  {\large
  \begin{tabular}{@{}ll@{}}
    \toprule
    \textbf{Dimension} & \textbf{IOI Circuit} \\
    \midrule
    Construct    & Partially confirmed \\
    Measurement  & \textcolor{disconfirmed}{\textbf{2 Disconfirmed}} \\
    Internal     & Partially confirmed \\
    External     & \textcolor{disconfirmed}{\textbf{1 Disconfirmed}} \\
    Interpretive & Partially confirmed \\
    \midrule
    \textbf{Verdict} & \textbf{Causally Suggestive} \\
    \bottomrule
  \end{tabular}}
  \end{center}

  \vspace{5mm}
  The IOI circuit has strong causal evidence. Method-conditional
  faithfulness (87\% under mean ablation, below 0\% to over
  100\% under others) caps the verdict at the external
  dimension.
}
}

\column{0.33}

\block{The Evidence Card}{
  A structured reporting unit that makes the gap visible.

  \vspace{4mm}
  \renewcommand{\arraystretch}{1.25}
  \begin{tabular}{@{}p{0.22\linewidth}p{0.70\linewidth}@{}}
    \toprule
    \textbf{Claim} & 26-head circuit implements IOI in GPT-2 Small \\
    \textbf{Claimed} & Triangulated \\
    \textbf{Supported} & \textcolor{tierCausal}{\textbf{Causally Suggestive}} \\
    \textbf{Caps it} & Faithfulness is method-conditional: 87\% mean ablation,
      $<$50\% under others \\
    \textbf{To close} & Cross-method agreement, specificity test,
      double dissociation \\
    \bottomrule
  \end{tabular}

  \vspace{6mm}
  Three questions for any mechanistic claim:

  \vspace{3mm}
  \begin{enumerate}[leftmargin=*, itemsep=3mm]
    \item Is the description mode declared?
    \item Is the claimed tier supported by the evidence?
    \item Do the two match?
  \end{enumerate}

  \vspace{4mm}
  Most published claims answer (1) implicitly, skip (2),
  and assume (3).
}

\end{columns}

% ════════════════════════════════════════════════════════
% ROW 2 — THE AUDIT + INDUCTION HEADS CASE
% ════════════════════════════════════════════════════════
\begin{columns}
\column{0.66}

{
\colorlet{blocktitlebgcolor}{colorOne}
\block{Sixteen Claims from Fifteen Papers}{
  {\small
  \renewcommand{\arraystretch}{1.15}
  \begin{tabular}{@{}>{\raggedright\arraybackslash}p{3.6cm}
                    >{\raggedright\arraybackslash}p{7.0cm}
                    >{\raggedright\arraybackslash}p{5.5cm}@{}}
    \toprule
    \textbf{Verdict} & \textbf{Claim} & \textbf{Scope} \\
    \midrule
    \textcolor{tierTriangulated}{\textbf{Triangulated}}
      & Induction heads: token copying
      & 1L--70B, random sequences \\
    \midrule
    \textcolor{tierSupported}{\textbf{Mech.\ Supported}}
      & Fourier mechanism (modular addition)
      & 1L ReLU, mod 113 \\
      & Greater-Than circuit
      & GPT-2 Small \\
      & Copy Suppression (L10H7)
      & GPT-2 Small \\
      & Superposition (toy networks)
      & 4 ReLU nets \\
      & Steering vector (refusal)
      & 13 models, 5 families \\
      & Global Workspace / J-space
      & 4 Claude models \\
      & Successor Heads
      & Pythia-1.4B \\
    \midrule
    \textcolor{tierCausal}{\textbf{Causally Suggestive}}
      & Othello World Model
      & Othello-GPT \\
      & Docstring Circuit
      & 4L attn-only \\
      & Gender Bias Circuits
      & GPT-2 + 5 families \\
      & IOI Circuit
      & GPT-2 Small \\
    \midrule
    \textcolor{tierProposed}{\textbf{Proposed}}
      & Probing Classifiers
      & review \\
    \midrule
    \textcolor{tierDisconfirmed}{\textbf{Disconfirmed}}
      & Knowledge Neurons
      & BERT-base \\
      & SAE features as canonical units
      & Pythia \\
      & Induction heads: general ICL
      & $>$1B \\
    \bottomrule
  \end{tabular}
  }

  \vspace{4mm}
  No claim reaches \textbf{Validated}. The highest tier
  achieved is Triangulated, by one claim.
  Same paper appears at both ends: induction heads for token
  copying (Triangulated) and for general ICL (Disconfirmed).
}
}

\column{0.34}

{
\colorlet{blocktitlebgcolor}{tierTriangulated}
\block{Same Paper, Two Verdicts}{
  Olsson et al.\ (2022) advance two claims of different reach.
  The evidence has since separated them.

  \vspace{4mm}
  \renewcommand{\arraystretch}{1.2}
  \begin{tabular}{@{}lcc@{}}
    \toprule
    & \textbf{Copying} & \textbf{General ICL} \\
    \midrule
    Confirmed          & 14 & 6 \\
    Partially conf.    & 15 & 14 \\
    Inconclusive       & 0  & 2 \\
    Untested           & 6  & 8 \\
    \textcolor{disconfirmed}{\textbf{Disconfirmed}} & \textbf{0} & \textbf{5} \\
    \midrule
    \textbf{Verdict} & \textcolor{tierTriangulated}{\textbf{Triang.}}
      & \textcolor{tierDisconfirmed}{\textbf{Disconf.}} \\
    \bottomrule
  \end{tabular}

  \vspace{5mm}
  \textbf{What flips:}
  \begin{itemize}[leftmargin=*, itemsep=2mm]
    \item C4 Discriminant: PC $\to$ D
    \item I1 Necessity: C $\to$ I
    \item E4 Cross-model: C $\to$ D
    \item V2 Level-evidence match: C $\to$ D
  \end{itemize}

  \vspace{4mm}
  The narrow mechanism replicates to 70B. The broad claim
  fails at scale---on evidence the authors named as their
  own refutation condition.
}
}

\end{columns}

% ════════════════════════════════════════════════════════
% ROW 3 — PROVENANCE + GAPS + TAKEAWAY
% ════════════════════════════════════════════════════════
\begin{columns}
\column{0.33}

\block{Where the Criteria Come From}{
  The criteria are imported, not invented. Each maps to
  an established standard:

  \vspace{4mm}
  \renewcommand{\arraystretch}{1.25}
  \begin{tabular}{@{}ll@{}}
    \toprule
    \textbf{Source} & \textbf{Criteria} \\
    \midrule
    Popper           & C1 Falsifiability \\
    Cronbach \& Meehl & C3, C4 Convergent / \\
                     & \quad Discriminant \\
    Campbell         & I1--I6 Internal \\
    Rosenbaum        & I7, I8 Confound \\
                     & \quad control / sensitivity \\
    Craver           & V1, V2 Level \\
                     & \quad declaration / match \\
    Benzer           & C6 Complementation \\
    \bottomrule
  \end{tabular}

  \vspace{5mm}
  What is new: the object being assessed (a causal
  mechanistic claim, not a test score or treatment effect),
  the interpretive criteria (V1--V5), and the treatment of a
  claim as a family of readings at different strengths.
}

\column{0.34}

{
\colorlet{blocktitlebgcolor}{tierCausal}
\block{What Caps the Middle Tier}{
  Every capped verdict names the criterion that caps it.

  \vspace{5mm}
  \begin{itemize}[leftmargin=*, itemsep=4mm]
    \item \textbf{IOI circuit} (E1): faithfulness is
      method-conditional. Same quantity, below 0\% to
      over 100\% across six choices.

    \item \textbf{SAE features} (C4): seed-dependent,
      no better than raw neurons at causal localization.

    \item \textbf{Othello} (V2): ``world model'' is an
      interpretive label on evidence supporting
      ``linearly decodable.''

    \item \textbf{Gender bias} (C4): presupposes bias
      separates from competence. Never tested.

    \item \textbf{Steering} (V2): extraction contrast
      is a harmfulness contrast, later separated from
      refusal by Zhao et al.
  \end{itemize}

  \vspace{4mm}
  Each gap is a missing \textbf{kind} of evidence, not a
  missing \textbf{amount}.
}
}

\column{0.33}

{
\colorlet{blocktitlebgcolor}{colorOne}
\block{Takeaway}{
  {\Large\bfseries Report the tier your evidence supports,
  and name the criterion that caps it.}

  \vspace{5mm}
  \begin{itemize}[leftmargin=*, itemsep=3mm]
    \item \textbf{36} criteria across \textbf{5} validity types
    \item \textbf{16} claims audited from \textbf{15} papers
    \item \textbf{1} Triangulated, \textbf{7} Supported,
      \textbf{4} Suggestive, \textbf{1} Proposed,
      \textbf{3} Disconfirmed
    \item \textbf{0} Validated
    \item Same paper at both ends of the scale
    \item The gap is always a missing \emph{kind} of evidence,
      not a missing amount
  \end{itemize}

  \vspace{6mm}
  \begin{center}
  \texttt{elliot@elliottower.ai}

  \vspace{3mm}
  {\small Framework, criteria, and all sixteen evidence cards:}

  \texttt{mechanisticresearch.org}
  \end{center}
}
}

\end{columns}

\end{document}
